% !TeX root = _main.tex
% ROLE 1: muc 1--4, van de, pham vi, he thong, nguon luc va AS-IS.
% Chi sua role cua minh va figures/rX_*. Khong sua tay bang do Python sinh.
% Nguoi ghep bai quan ly _main.tex va _content.tex.
% Macro Role*Supplement cuoi file giu ban nhap mo rong; KHONG duoc goi trong PDF.

\section{Bài toán, mục tiêu và phạm vi nghiên cứu}
\label{r1:scope}
\subsection{Bối cảnh và vấn đề cần giải quyết}
Hà Nội có 453 cửa hàng bán lẻ xăng dầu và nhu cầu tiêu thụ khoảng
110\,000~m$^3$ xăng/tháng theo số liệu được công bố tháng 3/2026~\cite{sct2026}.
Các số liệu này cung cấp quy mô tham chiếu, không mô tả lượng khách của một
cửa hàng cụ thể. Nhóm xét một cửa hàng giả định có 8 vòi bơm, phục vụ trong
khung 06:00--22:00, với lượng xe thay đổi theo giờ.

\textbf{Vấn đề dịch vụ.} Khách có thể phải chờ lâu hoặc bỏ đi vào giờ cao
điểm, trong khi một số giờ khác còn dư công suất. Báo cáo cần trả lời liệu
việc giữ số nhân viên cố định và thời gian chiếm trụ có góp phần tạo hàng
chờ hay không. Đây là giả thuyết được kiểm tra bằng mô phỏng; chưa phải
kết luận về một cửa hàng đã khảo sát.

\textbf{Ảnh hưởng.} Khách mất thời gian và có thể không mua được xăng;
nhân viên chịu áp lực không đều trong ngày; cửa hàng phải cân đối mức phục
vụ với chi phí nhân lực. Vì vậy, chỉ nhìn thời gian chờ trung bình cả ngày
hoặc tăng nhân viên cho mọi giờ đều chưa đủ căn cứ.

\subsection{Câu hỏi phân tích và mạch lập luận}
\textbf{Câu hỏi chính.} Có thể giảm tỷ lệ chờ quá 5 phút và tỷ lệ bỏ đi ở
khung 17:00--19:00 bằng cách điều chỉnh số kênh phục vụ theo giờ hay không,
và khi nào nên thử thêm cải tiến thanh toán?
\begin{table}[H]\centering\small
\caption{Mạch phân tích của báo cáo}\label{r1:problems}
\begin{tabularx}{\linewidth}{@{}p{28mm}Y Y@{}}
\toprule
Bước & Câu hỏi cần trả lời & Bằng chứng / đầu ra \\
\midrule
Nêu vấn đề & Chờ ở đâu, lúc nào, ảnh hưởng ai? & AS-IS, mốc đo, chỉ số theo giờ và tỷ lệ bỏ đi. \\
Phân tích nguyên nhân & Nhu cầu có vượt công suất? Trụ bị chiếm bởi thao tác nào? & $\lambda_t$, $c_t$, thời gian phục vụ; so sánh giả thuyết với mô hình. \\
Đề xuất giải pháp & Thay đổi gì để tác động đúng nguyên nhân? & A1 điều chỉnh công suất theo giờ; A2 thử giảm thời gian thanh toán tại trụ. \\
Đánh giá & Khách chờ ít hơn không? Chi phí và điều kiện thực hiện là gì? & Cùng dòng xe cho A0/A1/A2; độ nhạy, rủi ro và kế hoạch thử nghiệm. \\
\bottomrule
\end{tabularx}
\end{table}
Bài toán chính là hàng đợi và phân bổ công suất. Định vị cửa hàng, học máy
và chiến lược dài hạn được giữ ở bản nháp mở rộng, chưa dùng làm căn cứ cho
khuyến nghị của báo cáo này.

\subsection{Ranh giới hệ thống và ranh giới phép đo}
\begin{table}[htbp]\centering\small
\caption{Phạm vi dịch vụ và phạm vi mô hình định lượng}\label{r1:boundary}
\begin{tabularx}{\linewidth}{@{}p{32mm}Y@{}}
\toprule
Thành phần & Mô tả \\
\midrule
Khách hàng & Hai loại xe theo giả định: xe máy (97,3\% lượt) và ô tô (2,7\% lượt). \\
Tác nhân cung cấp & Nhân viên bơm, cửa hàng trưởng, doanh nghiệp đầu mối, ngân hàng/ví điện tử, cơ quan thuế. \\
Trigger dịch vụ & Xe vào cửa hàng và nhìn thấy hàng chờ tại các trụ bơm. \\
Outcome dịch vụ & Xe được bơm đủ, thanh toán xong, có hóa đơn điện tử và rời trụ. \\
Đầu vào mô hình & Thời điểm xe vào hàng chờ $a_i$ (hoặc bỏ đi nếu hàng quá dài). \\
Đầu ra mô hình & Thời điểm rời trụ $d_i$; $W_q=b_i-a_i$ với $b_i$ là lúc nhân viên bắt đầu phục vụ. \\
Ngoài mô hình & Chất lượng xăng, nhập hàng và tồn kho bể, an toàn phòng cháy, dịch vụ phụ (rửa xe, cửa hàng tiện lợi). \\
\bottomrule
\end{tabularx}
\end{table}
Cửa hàng có thể mở 24 giờ, nhưng phân tích chỉ xét 06:00--22:00, nơi tập
sinh toàn bộ lượt xe của kịch bản; nhu cầu ngoài giờ chưa được xét. Mỗi ``kênh'' trong mô hình là một vòi bơm có
nhân viên phục vụ; số kênh không vượt quá 8 vòi vật lý.

\subsection{Mục tiêu và tiêu chí đánh giá}
Mục tiêu là cải thiện khả năng phục vụ giờ cao điểm với chi phí nhân lực
hợp lý. Nhóm đề xuất các ngưỡng để so sánh phương án: thời gian chờ trung
bình không quá 2 phút, tỷ lệ chờ quá 5 phút không quá 5\%, tỷ lệ bỏ đi không
quá 1\%. Đây là mục tiêu nghiên cứu, cần thống nhất lại với cửa hàng nếu
triển khai. Báo cáo đánh giá riêng cả ngày và khung 17:00--19:00, đồng thời
so sánh giờ công, chi phí và mức chịu tải khi nhu cầu tăng.

Không đặt trước một tỷ lệ cắt giảm nhân lực phải đạt: phương án chỉ được
đề xuất nếu giữ chất lượng phục vụ và có thể chuyển thành lịch ca thực tế.

\section{Mô hình hệ thống dịch vụ}
\label{r1:system}
Báo cáo dùng mô hình hệ thống dịch vụ làm khung chủ đạo để liên kết tác nhân,
nguồn lực, hoạt động, thông tin, quy tắc và kết quả~\cite{rubric}.
\begin{table}[htbp]\centering\small
\caption{Các thành phần của hệ thống và liên hệ với bài toán chờ}\label{r1:components}
\begin{tabularx}{\linewidth}{@{}p{24mm}Y Y@{}}
\toprule
Thành phần & Biểu hiện tại cửa hàng & Liên hệ với phân tích \\
\midrule
Con người & Khách xe máy/ô tô, nhân viên bơm, cửa hàng trưởng & Số kênh có nhân viên $c_t$; tỷ trọng loại xe. \\
Hoạt động & Mời xe vào trụ, bơm, thu tiền, xuất hóa đơn & Thời gian phục vụ $S$ theo loại xe. \\
Thông tin & Giá niêm yết, số tiền/số lít, hóa đơn điện tử & Dữ liệu từng lần bán dùng để dự báo $\lambda_t$. \\
Công nghệ & Cột bơm, máy POS, QR/ví, ứng dụng trả trước & Số hóa giảm thời gian thanh toán (A2). \\
Quy tắc & FCFS tại mỗi trụ, hóa đơn từng lần bán, giờ mở cửa & Ràng buộc $c_t\le8$; mục tiêu thời gian chờ. \\
Kết quả & Xe được phục vụ, khách quay lại & $W_q$, $P(W_q{>}5)$, tỷ lệ bỏ đi, chi phí. \\
\bottomrule
\end{tabularx}
\end{table}

\section{Bản đồ nguồn lực và cơ chế đồng tạo giá trị}
\label{r1:resources}
\begin{figure}[htbp]\centering
\resizebox{0.95\linewidth}{!}{%
\begin{tikzpicture}[node distance=10mm,
  hub/.style={rectangle,rounded corners=4pt,draw,very thick,fill=blue!6,align=center,font=\small,minimum width=34mm,minimum height=12mm},
  act/.style={rectangle,rounded corners=3pt,draw,thick,fill=gray!8,align=center,font=\scriptsize,minimum width=27mm,minimum height=9mm},
  res/.style={font=\scriptsize\itshape,align=center,fill=white,inner sep=1pt}]
\node[hub] (s) at (0,0) {Cửa hàng xăng dầu\\(8 vòi, nhân viên bơm)};
\node[act] (c) at (-6.2,2.4) {Khách xe máy / ô tô\\(theo thế hệ)};
\node[act] (b) at (6.2,2.4) {Doanh nghiệp đầu mối\\(nguồn hàng, chiết khấu)};
\node[act] (p) at (-6.2,-2.4) {Ngân hàng / ví điện tử\\(QR, trả trước)};
\node[act] (t) at (6.2,-2.4) {Cơ quan thuế,\\Sở Công Thương};
\node[act] (e) at (0,-3.6) {Đối tác trạm sạc / đổi pin};
\draw[<->,thick] (c) -- node[res]{tiền, thời gian,\\dữ liệu mua hàng} (s);
\draw[<->,thick] (b) -- node[res]{xăng, giá bán lẻ\\$\leftrightarrow$ sản lượng} (s);
\draw[<->,thick] (p) -- node[res]{hạ tầng thanh toán\\$\leftrightarrow$ phí giao dịch} (s);
\draw[<->,thick] (t) -- node[res]{quy định, hóa đơn\\từng lần bán} (s);
\draw[<->,thick,dashed] (e) -- node[res]{mặt bằng $\leftrightarrow$ khách xe điện} (s);
\end{tikzpicture}}
\caption{Bản đồ nguồn lực của cửa hàng xăng (nét đứt: quan hệ tương lai khi chuyển dịch xe điện).}
\label{r1:resource-map}
\source{nhóm vẽ bằng TikZ dựa trên các tác nhân trong Bảng~\ref{r1:components}; Claude Code hỗ trợ dựng mã hình}
\end{figure}
Khách hàng góp thời gian chờ, tiền và dữ liệu mua hàng; nhân viên góp thao tác
phục vụ; đầu mối cung cấp nguồn hàng và chiết khấu; ngân hàng/ví cung cấp
thanh toán số; cơ quan thuế đặt quy tắc hóa đơn. Giá trị (``vào là được bơm
ngay'') chỉ xuất hiện khi các nguồn lực này khớp với nhau theo giờ: thiếu nhân
viên giờ tan tầm hoặc thao tác thanh toán chậm đều chuyển chi phí sang thời
gian chờ của khách. Hai cơ chế này tương ứng với hai biến quyết định $c_t$ và
$E[S]$ của mô hình.

\begin{table}[htbp]\centering\small
\caption{Lợi ích các bên và điều kiện chấp nhận giải pháp}\label{r1:stakeholders}
\begin{tabularx}{\linewidth}{@{}p{25mm}Y Y Y@{}}
\toprule
Bên liên quan & Lợi ích hiện tại & Kỳ vọng dự kiến & Điều kiện cải tiến \\
\midrule
Khách dùng QR & Thanh toán thuận tiện & Chờ ít hơn, giao dịch chính xác & Đo thời gian và lỗi giao dịch khi thử A2 \\
Khách dùng tiền mặt & Nhân viên hỗ trợ, trả tiền mặt & Vẫn được phục vụ khi không dùng ứng dụng & Giữ kênh tiền mặt, không ép số hóa \\
Nhân viên bơm & Lương 8--12 triệu đồng/tháng~\cite{wage} & Ca ổn định, không quá tải giờ cao điểm & Lịch theo giờ cần báo trước, có ca linh hoạt có thù lao \\
Đầu mối cung ứng & Sản lượng ổn định & Cung ứng đúng kế hoạch & Phân biệt thiếu hàng với thiếu công suất phục vụ \\
Cơ quan quản lý & Giao dịch minh bạch & Thông tin đầy đủ, có thể đối soát & Không bỏ bước kiểm tra để rút ngắn thời gian \\
\bottomrule
\end{tabularx}
\end{table}

\section{Quy trình hiện tại và chẩn đoán điểm nghẽn}
\label{r1:asis}
\subsection{Hành trình khách hàng và blueprint AS-IS}
Hành trình gồm: thấy cửa hàng $\rightarrow$ quyết định vào hay đi tiếp (nhìn độ
dài hàng) $\rightarrow$ chờ $\rightarrow$ mở nắp bình, báo số tiền $\rightarrow$
bơm $\rightarrow$ thanh toán, nhận hóa đơn $\rightarrow$ rời trụ.
Hai điểm đau chính là \emph{không đoán được phải chờ bao lâu} khi nhìn hàng
và \emph{thao tác thanh toán bằng tiền mặt, trả lại tiền thừa, xuất hóa đơn}
kéo dài thời gian chiếm trụ.
\begin{table}[htbp]\centering\small
\caption{Các lớp của blueprint tại điểm chờ}\label{r1:blueprint}
\begin{tabularx}{\linewidth}{@{}p{34mm}Y@{}}
\toprule
Lớp & Nội dung \\
\midrule
Bằng chứng hữu hình & Biển giá, đồng hồ cột bơm, hàng xe đang chờ, hóa đơn điện tử. \\
Hành động khách & Chọn trụ, xếp hàng, mở nắp bình, báo số tiền, thanh toán. \\
\multicolumn{2}{@{}c@{}}{\footnotesize\textit{--- đường tương tác ---}} \\
Tuyến trước & Nhân viên mời xe, bơm, thu tiền, trả tiền thừa. \\
\multicolumn{2}{@{}c@{}}{\footnotesize\textit{--- đường nhìn thấy ---}} \\
Tuyến sau & Máy POS ghi giao dịch, xuất hóa đơn điện tử từng lần bán. \\
Quy trình hỗ trợ & Lịch ca, nhập hàng, đối soát tiền mặt cuối ca, kết nối cơ quan thuế. \\
\bottomrule
\end{tabularx}
\end{table}

\subsection{Sơ đồ quy trình AS-IS (BPMN rút gọn)}
\begin{figure}[htbp]\centering
\resizebox{\linewidth}{!}{%
\begin{tikzpicture}
\def\W{18.4}
\foreach \y/\name in {0/Khách hàng,-2.4/Nhân viên bơm,-4.8/Hệ thống POS / HĐĐT}{
  \draw (0,\y) rectangle (\W,\y-2.4);
  \draw (0.8,\y) -- (0.8,\y-2.4);
  \node[rotate=90,font=\scriptsize\bfseries] at (0.4,\y-1.2) {\name};}
\node[ev] (st) at (1.6,-1.2) {};
\node[task] (t1) at (3.4,-1.2) {Xe vào\\cửa hàng};
\node[gw] (g1) at (5.6,-1.2) {Hàng chờ\\$\ge$ 20 xe?};
\node[evend] (eb) at (7.3,-0.45) {};
\node[font=\scriptsize,anchor=west] at (7.7,-0.45) {bỏ đi (balk)};
\node[task] (t2) at (8.4,-1.6) {Xếp hàng\\chờ tới lượt};
\node[task] (t3) at (10.6,-3.6) {Mời xe vào trụ,\\hỏi số tiền/số lít};
\node[task] (t4) at (13.0,-3.6) {Bơm xăng};
\node[task] (t5) at (13.0,-1.2) {Thanh toán\\tiền mặt / QR};
\node[task] (t6) at (15.6,-6.0) {Xuất hóa đơn\\điện tử từng lần};
\node[task] (t7) at (15.6,-1.2) {Nhận hóa đơn,\\rời trụ};
\node[evend] (en) at (17.7,-1.2) {};
\draw[flow] (st) -- (t1);
\draw[flow] (t1) -- (g1);
\draw[flow] (g1.north) |- node[pos=0.25,left,font=\scriptsize]{có} (eb);
\draw[flow] (g1.east) -- node[above,font=\scriptsize]{không} ++(0.6,0) |- (t2);
\draw[flow] (t2) |- (t3);
\draw[flow] (t3) -- (t4);
\draw[flow] (t4) -- (t5);
\draw[flow] (t5.east) -- ++(0.35,0) |- (t6.west);
\draw[flow] (t6.north) -- (t7.south);
\draw[flow] (t7) -- (en);
\node[tmark,anchor=west] at (9.5,-1.6) {$a_i$: vào hàng};
\node[tmark] at (10.6,-4.35) {$b_i$: bắt đầu phục vụ};
\node[tmark] at (15.6,-0.35) {$d_i$: rời trụ};
\end{tikzpicture}}
\caption{Quy trình AS-IS tại một vòi bơm; nhãn đỏ là các mốc thời gian dùng để đo $W_q=b_i-a_i$ và $S=d_i-b_i$.}
\label{r1:bpmn}
\source{nhóm vẽ bằng TikZ theo ký pháp BPMN rút gọn (sự kiện, tác vụ, cổng XOR, làn); Claude Code hỗ trợ dựng mã hình}
\end{figure}
Thời gian phục vụ $S$ gồm cả bơm, thanh toán và xuất hóa đơn: trụ chỉ được giải
phóng khi khách rời đi. Nếu giảm được thao tác tại trụ, A2 có thể rút ngắn
thời gian chiếm trụ dù tốc độ bơm không đổi.
\todo{1}{Đối chiếu sơ đồ với quan sát tại một cửa hàng thật (bấm giờ 30--50 lượt
mỗi loại xe), sửa nhánh ngoại lệ nếu có (hết loại xăng, lỗi POS).}

\subsection{Từ điểm nghẽn đến biến đo}
\begin{table}[htbp]\centering\small
\caption{Giả thuyết chẩn đoán và biến kiểm tra}\label{r1:diagnosis}
\begin{tabularx}{\linewidth}{@{}Y Y Y@{}}
\toprule
Điểm trên quy trình & Cơ chế giả định & Biến kiểm tra \\
\midrule
Hàng dài giờ tan tầm & Lịch nhân viên cố định trong khi lượt xe dao động theo giờ & $\lambda_t$, $c_t$, $\rho_t$, $W_q$ theo giờ \\
Khách bỏ đi & Thấy hàng dài nên sang cửa hàng khác & Tỷ lệ bỏ đi (balking) \\
Chiếm trụ lâu & Tiền mặt, tiền thừa, hóa đơn từng lần bán & $E[S]$ theo loại xe; CV \\
Ngày có sự kiện & Lượng xe tăng ngoài cấu hình cơ sở & $\lambda_t$, hàng chờ và tỷ lệ bỏ đi \\
\bottomrule
\end{tabularx}
\end{table}


% Ban nhap bo tro: chua dua vao PDF chinh.
\newcommand{\RoleOneSupplement}{%
\subsection{Môi trường chiến lược: SWOT và PESTEL}
\begin{table}[htbp]\centering\small
\caption{SWOT của cửa hàng đại diện}\label{r1:swot}
\begin{tabularx}{\linewidth}{@{}Y Y@{}}
\toprule
\textbf{Điểm mạnh (S)} & \textbf{Điểm yếu (W)} \\
Vị trí nội thành, lượng khách lớn; có dữ liệu hóa đơn từng lần bán để đo và dự báo nhu cầu. &
Lịch nhân viên cố định; thanh toán tiền mặt kéo dài thời gian chiếm trụ; tối đa 8 vòi. \\
\midrule
\textbf{Cơ hội (O)} & \textbf{Thách thức (T)} \\
Thanh toán QR/ứng dụng trả trước; dùng mặt bằng cho trạm sạc, đổi pin~\cite{vgx}; đối thủ đóng cửa khi chiết khấu thấp. &
Chuyển dịch sang xe điện và vùng phát thải thấp~\cite{ct20,lez2026}; đứt gãy nguồn cung~\cite{vtc2026}; biên lợi nhuận mỏng~\cite{chietkhau}. \\
\bottomrule
\end{tabularx}
\end{table}

\begin{table}[htbp]\centering\small
\caption{PESTEL: yếu tố dẫn dắt và kịch bản tương ứng ở Mục~\ref{r3:scenarios}}\label{r1:pestel}
\begin{tabularx}{\linewidth}{@{}p{21mm}Y p{17mm}@{}}
\toprule
Yếu tố & Biểu hiện có nguồn & Kịch bản \\
\midrule
Chính trị (P) & Chỉ thị 20/CT-TTg; Đề án vùng phát thải thấp trong Vành đai~1 từ 1/7/2026, cấm theo khung giờ tại khu vực thí điểm~\cite{ct20,lez2026} & KB-A \\
Kinh tế (E) & Giá điều chỉnh hằng tuần (kỳ 5/11/2026 giảm 2\,075 đồng/lít, kỳ 12/11 giảm 1\,310 đồng/lít~\cite{moit2026}); chiết khấu bán lẻ thấp~\cite{chietkhau} & KB-B, KB-C \\
Xã hội (S) & Hà Nội đăng ký mới hơn 255\,000 xe máy năm 2025~\cite{reg2025}; xe máy xăng vẫn chiếm 97\% trong Vành đai~1~\cite{xe97} & KB-B \\
Công nghệ (T) & Thanh toán QR, hóa đơn điện tử; mạng lưới đổi pin tại cửa hàng xăng~\cite{vgx} & KB-A \\
Môi trường (E) & Mục tiêu giảm phát thải giao thông đô thị & KB-A \\
Pháp lý (L) & Hóa đơn điện tử từng lần bán (NĐ 70/2025)~\cite{nd70}; lộ trình cấm có thể lùi tới 2028~\cite{tt2028} & Mọi KB \\
\bottomrule
\end{tabularx}
\end{table}
Các yếu tố trên cho thấy hai bất định lớn: \emph{chiều nhu cầu} (giảm do xe
điện hay tăng do đối thủ đóng cửa) và \emph{cú sốc ngắn hạn} (đứt gãy nguồn
cung, điều chỉnh giá). Đây là cơ sở để xây ba kịch bản A/B/C ở Mục~\ref{r3:scenarios}.
\todo{1}{Mở và đọc lại từng nguồn báo chí/văn bản trong \texttt{thebibliography}
trước khi nộp; nếu nội dung khác số liệu trích, báo người ghép bài sửa.}

}
